% \documentclass{standalone} is particular class that non have a size, but the size is define from
% the document itself
\documentclass[tikz,margin=0pt, convert={outfile=SFOM_10-60.png}]{standalone}
% \usepackage{tikz} is the package that contain the graphics command as \draw etc ...
% The package is describe in this document: https://www.sharelatex.com/learn/TikZ_package
\usepackage{tikz}
% If the graduate dials is for analogical gauges period 50-70 years, this font is ok
% OCR is open-source package deposited in the CTAN repository: https://www.ctan.org/pkg/ocr-b
% Opentype can be used on Linux and Mac machines
% (I do not know if it's possible in Windows, but I think of it),
% you can always download it from the CTAN site at this address:
% https://www.ctan.org/tex-archive/fonts/ocr-b-outline
\usepackage{ocr}
% Defines the font size to use, the first number defines the font size,
% the second number the dimensional family from which the font is taken.
% It can be assumed that there are differences between a dimensional family 10 and a 12 dimensional font.
\tikzset{font={\fontsize{7pt}{10}\selectfont}}
\begin{document}
% The link OCR package documentation is this: http://ctan.mirror.garr.it/mirrors/CTAN/macros/latex/contrib/ocr-latex/ocr.pdf
% \ocrfamily - Normal font family
% \ocrnegfamily - Family negative fonts
{\ocrfamily
    % Xscale defines the scale of the drawing with respect to the document we write.
    % It is imperative to make it compatible for the size of the graduated dial that we are designing.
    % Its determination requires a number of tests to optimize the font dimension with the units of measurement.
    \begin{tikzpicture}[xscale = 0.0480]
      %5° Rays
      \foreach \a in {-180, -179,...,180}
	\draw[line width=0.3mm] (\a,0) -- (\a,2mm);
      %10° Rays
      \foreach \a in {-10, -5,...,60}
	\draw[line width=0.3mm] (\a,0) -- (\a,3mm);
      %30° Rays
      \foreach \a in {-10, 0,...,60}
	\draw[line width=0.3mm] (\a,0) -- (\a,4mm); 
      %labels  
      \draw (-10,0.7) node[rotate=90] { 10};
      \draw (0,0.7) node[rotate=90] { 0};
      \draw (10,0.7) node[rotate=90] {-10};
      \draw (20,0.7) node[rotate=90] {-20};
      \draw (30,0.7) node[rotate=90] {-30};
      \draw (40,0.7) node[rotate=90] {-40};
      \draw (50,0.7) node[rotate=90] {-50};
      \draw (60,0.7) node[rotate=90] {-60};
    \end{tikzpicture}
}
\end{document} 
